\chapter{Composite Bosons, Spin Representations, and Dirac Doubling}

\section{Exchange Symmetry and Composite Bosons}

The opening question is a genuine paradox about composite particles. If two fermions make a boson, can two such bosons occupy the same state without violating fermionic exclusion? The answer begins with the exchange rule itself.

Suppose two identical particles are described by a wavefunction of two arguments. The symbols \(x\) and \(y\) label the two particle slots; they do not distinguish two species. For bosons the wavefunction is symmetric, while for fermions it is antisymmetric:
\begin{equation}
\Psi_B(x,y)=\Psi_B(y,x),
\qquad
\Psi_F(x,y)=-\Psi_F(y,x).
\end{equation}
The symmetry says that interchanging the two indistinguishable particles does not produce a new physical state.

Two bosons can certainly occupy the same one-particle state. If that state has wavefunction \(\varphi\), the two-particle product is
\begin{equation}
\Psi(x,y)=\varphi(x)\varphi(y).
\end{equation}
Wavefunctions are ordinary complex-valued functions, not operators or matrices, so their values commute:
\begin{equation}
\varphi(x)\varphi(y)=\varphi(y)\varphi(x).
\end{equation}
The product is therefore symmetric. The same construction cannot give a nonzero two-fermion state, because antisymmetry would require
\begin{equation}
\varphi(x)\varphi(y)=-\varphi(y)\varphi(x),
\end{equation}
while the two sides without the minus sign are equal.

More generally, begin with any two-particle wavefunction \(F(x,y)\), with no assumed exchange symmetry. Its symmetric and antisymmetric parts are, apart from normalization,
\begin{align}
F_{\mathrm{sym}}(x,y)&=F(x,y)+F(y,x),\\
F_{\mathrm{asym}}(x,y)&=F(x,y)-F(y,x).
\end{align}
The exchange check is explicit:
\begin{align}
F_{\mathrm{sym}}(y,x)
&=F(y,x)+F(x,y)
=F_{\mathrm{sym}}(x,y),\\
F_{\mathrm{asym}}(y,x)
&=F(y,x)-F(x,y)
=-F_{\mathrm{asym}}(x,y).
\end{align}
Adding the exchanged function symmetrizes it; subtracting the exchanged function antisymmetrizes it.

A pair of hydrogen atoms provides the more interesting case. There are four constituent fermions: two identical electrons and two identical protons. Let \(x_1,x_2\) denote the electron variables and \(y_1,y_2\) the proton variables. The four-body wavefunction must be antisymmetric separately under exchange of the electrons and under exchange of the protons:
\begin{align}
\Psi(x_2,x_1;y_1,y_2)
&=-\Psi(x_1,x_2;y_1,y_2),\\
\Psi(x_1,x_2;y_2,y_1)
&=-\Psi(x_1,x_2;y_1,y_2).
\end{align}

To motivate two distinct one-atom wavefunctions, imagine one hydrogen atom localized near the origin, with wavefunction \(\psi(x,y)\), and another localized elsewhere, with wavefunction \(\phi(x,y)\). The second may simply be a translated copy of the first, obtained by translating both its electron and proton arguments, but it is still a distinct wavefunction in this setup. A naive product is
\begin{equation}
\psi(x_1,y_1)\phi(x_2,y_2).
\end{equation}
It does not yet have either required fermionic antisymmetry.

First antisymmetrize the electron variables:
\begin{equation}
\Psi_x=
\psi(x_1,y_1)\phi(x_2,y_2)
-\psi(x_2,y_1)\phi(x_1,y_2).
\end{equation}
Under \(x_1\leftrightarrow x_2\), the two terms exchange places and the whole expression changes sign. Now antisymmetrize this result in the proton variables. The four-term ordering is fixed by the two successive antisymmetrizations.\footnote{This construction also fixes the sign pattern \(+,-,-,+\) through the variable-label corrections made during the spoken calculation.}
\begin{align}
\Psi_{2H}(x_1,x_2;y_1,y_2)
&=\Psi_x(x_1,x_2;y_1,y_2)
-\Psi_x(x_1,x_2;y_2,y_1)\\
&=\psi(x_1,y_1)\phi(x_2,y_2)
-\psi(x_2,y_1)\phi(x_1,y_2)\nonumber\\
&\quad
-\psi(x_1,y_2)\phi(x_2,y_1)
+\psi(x_2,y_2)\phi(x_1,y_1).
\label{eq:two-hydrogen-antisymmetrized}
\end{align}
An electron exchange interchanges the first two terms and the last two terms, with an overall sign reversal. A proton exchange does the same. Thus the state has exactly the two separate antisymmetries required of the constituent fermions.

\begin{classroomqa}
\lecturetimestamp{00:11:11}
\audiencequestion{If hydrogen atoms are bosons, what prevents two of them from being placed directly on top of each other?}
\lecturerresponse{Bosonic exchange statistics do not remove the forces between the atoms. Pauli exclusion among the constituents, together with the other interatomic forces, can make spatial overlap extremely costly. That is different from asking whether the atoms may occupy the same delocalized momentum state. A momentum eigenstate is spread throughout space, so common momentum does not mean ramming the atoms into the same point. Even hard-core objects like billiard balls could obey bosonic exchange statistics: their hard cores would prevent spatial overlap, while bosonic statistics would still allow common momentum occupation.}
\end{classroomqa}

\begin{classroomqa}
\lecturetimestamp{00:12:28}
\audiencequestion{Would putting two hydrogen atoms in the same momentum eigenstate put their constituent fermions in the same state?}
\lecturerresponse{No. The four-body state remains antisymmetric under exchange of the two electron variables and separately antisymmetric under exchange of the two proton variables. Only the simultaneous exchange of both constituents---the exchange of the two whole atoms---is symmetric.}
\end{classroomqa}

Indeed, exchanging the two atoms means making both substitutions
\begin{equation}
x_1\leftrightarrow x_2,
\qquad
y_1\leftrightarrow y_2.
\end{equation}
The first substitution contributes one fermionic minus sign and the second contributes another:
\begin{equation}
\Psi_{2H}(x_2,x_1;y_2,y_1)
=(-1)(-1)\Psi_{2H}(x_1,x_2;y_1,y_2)
=\Psi_{2H}(x_1,x_2;y_1,y_2).
\end{equation}
The same result can be seen term by term in Eq.~\eqref{eq:two-hydrogen-antisymmetrized}. Under the whole-atom exchange,
\begin{equation}
\psi(x_1,y_1)\phi(x_2,y_2)
\longleftrightarrow
\psi(x_2,y_2)\phi(x_1,y_1),
\end{equation}
so the first and fourth terms exchange. Likewise,
\begin{equation}
-\psi(x_2,y_1)\phi(x_1,y_2)
\longleftrightarrow
-\psi(x_1,y_2)\phi(x_2,y_1),
\end{equation}
so the second and third terms exchange. The complete state is symmetric under interchange of the atoms. Two constituent minus signs give the composite a plus sign.

This exchange statement is logically independent of the energy required to overlap the atoms. Strong repulsion may make such overlap inaccessible, but it does not change the symmetric exchange law.

\begin{classroomqa}
\lecturetimestamp{00:14:42}
\audiencequestion{What does the constructed four-body state represent physically, and is it specifically a momentum-space state?}
\lecturerresponse{It is the two-atom wavefunction obtained from two one-atom wavefunctions by antisymmetrizing the electron labels and then the proton labels. It need not be written in momentum space. The variables may describe position, momentum, or another representation; the required constituent exchange symmetries are the same.}
\end{classroomqa}

Consequently, two atoms may occupy the same delocalized momentum wavefunction even though forcing their spatial distributions into strong overlap can cost a great deal of energy.

Now set the two one-atom wavefunctions equal, \(\phi=\psi\). Equation~\eqref{eq:two-hydrogen-antisymmetrized} becomes
\begin{align}
\Psi_{\mathrm{same}}
&=\psi(x_1,y_1)\psi(x_2,y_2)
-\psi(x_2,y_1)\psi(x_1,y_2)\nonumber\\
&\quad
-\psi(x_1,y_2)\psi(x_2,y_1)
+\psi(x_2,y_2)\psi(x_1,y_1).
\end{align}

Does this four-body state vanish? It does not. The first and fourth products are equal, and the second and third products are equal. Using the commutativity of the one-atom wavefunctions,
\begin{equation}
\Psi_{\mathrm{same}}
=
2\left[
\psi(x_1,y_1)\psi(x_2,y_2)
-\psi(x_2,y_1)\psi(x_1,y_2)
\right],
\label{eq:same-hydrogen-state}
\end{equation}
which is generally nonzero. It remains antisymmetric in the two electron variables and separately antisymmetric in the two proton variables.

There is no assignment here of one sharply located electron to each sharply located atom. An electron bound in an atom is spatially delocalized over its atomic wavefunction. Giving the two atoms the same atomic state therefore does not set \(x_1=x_2\), nor does it set \(y_1=y_2\).

The constituent exclusion is visible by testing coincidence directly. If the two modeled electron arguments coincide, \(x_1=x_2=x\), Eq.~\eqref{eq:same-hydrogen-state} gives
\begin{align}
\Psi_{\mathrm{same}}(x,x;y_1,y_2)
&=2\left[
\psi(x,y_1)\psi(x,y_2)
-\psi(x,y_1)\psi(x,y_2)
\right]\\
&=0.
\end{align}
The corresponding expression also vanishes when \(y_1=y_2\). Thus the same-state composite wavefunction is nonzero, while its amplitude at coincident identical-fermion arguments vanishes.\footnote{Spin and any other one-particle labels are suppressed in this model. Fermionic antisymmetry applies to the complete set of one-particle labels; the displayed cancellation concerns coincidence of the arguments retained here.}

\begin{classroomqa}
\lecturetimestamp{00:20:41}
\audiencequestion{Can the atoms be in the same state while their electrons and protons are not in the same constituent states?}
\lecturerresponse{Yes. Equality of the composite atomic wavefunctions does not remove the separate electron and proton antisymmetries. The complete four-body state can describe two atoms in one atomic state while assigning zero amplitude to coincident identical-fermion arguments.}
\end{classroomqa}

\begin{classroomqa}
\lecturetimestamp{00:21:12}
\audiencequestion{Did antisymmetrization project out the pieces in which the identical constituents coincide?}
\lecturerresponse{Yes. In the variables displayed here, those coincidence components cancel exactly. The precise statement is that the antisymmetric wavefunction vanishes when the two identical-fermion arguments are equal.}
\end{classroomqa}

The cancellation did not depend on first making the atomic wavefunctions equal. For the general state in Eq.~\eqref{eq:two-hydrogen-antisymmetrized},
\begin{align}
\Psi_{2H}(x,x;y_1,y_2)
&=\psi(x,y_1)\phi(x,y_2)
-\psi(x,y_1)\phi(x,y_2)\nonumber\\
&\quad
-\psi(x,y_2)\phi(x,y_1)
+\psi(x,y_2)\phi(x,y_1)\\
&=0.
\end{align}

\begin{classroomqa}
\lecturetimestamp{00:21:32}
\audiencequestion{Does this cancellation still hold when the atoms are in different states and the identical constituent variables coincide?}
\lecturerresponse{Yes. It follows from antisymmetry itself, not from setting \(\phi=\psi\). The same statement holds in another representation: after transforming to momentum space, the amplitude vanishes when the two identical-fermion momentum arguments coincide, provided the other suppressed one-particle labels coincide as well.}
\end{classroomqa}

\section{Spin as a Finite-Dimensional Quantum System}

The mathematics of spin recurs in isospin, color, and the other internal symmetries of particle physics. The analogy is formal: the same algebra appears, although the physical quantities need not have the same interpretation. The immediate task is to develop spin a little further and then return to the Dirac equation. This machinery will later be used for isospin, a concept associated historically with the 1930s.

An electron is a spin-\(\tfrac12\) particle, as are the proton and the muon. Its \(z\)-component of spin has two possible values. Before setting units, write
\begin{equation}
S_z=m\hbar.
\end{equation}
From now on \(\hbar=1\). The letter \(L\) is conventionally reserved for orbital angular momentum, so intrinsic spin will be denoted by \(S\). It obeys the same angular-momentum algebra:
\begin{equation}
[S_x,S_y]=iS_z,
\qquad
[S_y,S_z]=iS_x,
\qquad
[S_z,S_x]=iS_y,
\end{equation}
and
\begin{equation}
S_\pm=S_x\pm iS_y.
\end{equation}
The operator \(S_+\) raises the \(S_z\) eigenvalue and \(S_-\) lowers it.

\begin{classroomqa}
\lecturetimestamp{00:25:39}
\audiencequestion{Is the discussion specifically about intrinsic spin?}
\lecturerresponse{Intrinsic spin is the intended focus. The same angular-momentum algebra applies more generally, including to orbital angular momentum or the rotation of a macroscopic object.}
\end{classroomqa}

For spin one-half the only \(S_z\) eigenvalues are \(+\tfrac12\) and \(-\tfrac12\). Suppress everything else about the electron for the moment: forget its position and momentum and retain only the two spin states. Then
\begin{equation}
|\psi\rangle
=
\alpha|\uparrow_z\rangle+\beta|\downarrow_z\rangle
\quad\longleftrightarrow\quad
\begin{pmatrix}
\alpha\\
\beta
\end{pmatrix},
\end{equation}
where
\begin{equation}
|\alpha|^2=P(\uparrow_z),
\qquad
|\beta|^2=P(\downarrow_z),
\qquad
|\alpha|^2+|\beta|^2=1.
\end{equation}
This is a two-dimensional state space. It is not a claim that ordinary space has two dimensions.

\begin{classroomqa}
\lecturetimestamp{00:28:18}
\audiencequestion{Was physical space said to be two-dimensional?}
\lecturerresponse{No. Spin state space is two-dimensional. Position and momentum have temporarily been suppressed, as if the electron were nailed to a wall and only its spin could be measured.}
\end{classroomqa}

Operators on this state space are represented by \(2\times2\) matrices. It is convenient to choose the \(S_z\) basis so that \(S_z\) is diagonal, but diagonality is a basis choice rather than a physical requirement.

\begin{classroomqa}
\lecturetimestamp{00:29:46}
\audiencequestion{Must \(S_z\) be diagonal?}
\lecturerresponse{No. It is chosen diagonal for convenience. Other matrix representations are physically equivalent.}
\end{classroomqa}

\begin{classroomqa}
\lecturetimestamp{00:30:08}
\audiencequestion{Is \(S_z\) a matrix?}
\lecturerresponse{\(S_z\) is an observable and therefore an operator. In the chosen finite-dimensional basis, that operator is represented by a matrix acting on state vectors.}
\end{classroomqa}

A particularly useful representation is
\begin{equation}
S_i=\frac12\sigma_i,
\end{equation}
where
\begin{equation}
\sigma_x=
\begin{pmatrix}
0&1\\
1&0
\end{pmatrix},
\qquad
\sigma_y=
\begin{pmatrix}
0&-i\\
i&0
\end{pmatrix},
\qquad
\sigma_z=
\begin{pmatrix}
1&0\\
0&-1
\end{pmatrix}.
\end{equation}
The matrices without the factor of one-half are the Pauli matrices; the spin operators themselves include the factor. Direct multiplication verifies the commutation relations. The representation is not unique: rotating the axes or changing basis produces another equivalent triplet.

This is a two-component representation of spin-one-half angular momentum.

\subsection{Eigenstates along the coordinate axes}

It is useful to solve the eigenvalue equations component by component. Since the eigenvalues of \(\sigma_i\) are \(\pm1\), those of \(S_i=\sigma_i/2\) are \(\pm\tfrac12\).

For \(\sigma_z\), the two equations are
\begin{align}
\begin{pmatrix}
1&0\\
0&-1
\end{pmatrix}
\begin{pmatrix}
\alpha\\
\beta
\end{pmatrix}
&=
\begin{pmatrix}
\alpha\\
\beta
\end{pmatrix},
&
\begin{pmatrix}
1&0\\
0&-1
\end{pmatrix}
\begin{pmatrix}
\alpha\\
\beta
\end{pmatrix}
&=
-
\begin{pmatrix}
\alpha\\
\beta
\end{pmatrix}.
\end{align}

\begin{figure}[t]
\centering
\includegraphics[width=0.74\textwidth]{lecture_08_figure_02.png}
\caption{The \(+1\) eigenvalue equation for \(\sigma_z\) acting on a two-component spinor.}
\lectureframe{00:34:57}
\end{figure}

For eigenvalue \(+1\), row-by-row multiplication gives
\begin{equation}
\begin{pmatrix}
1\cdot\alpha+0\cdot\beta\\
0\cdot\alpha-\beta
\end{pmatrix}
=
\begin{pmatrix}
\alpha\\
-\beta
\end{pmatrix}
=
\begin{pmatrix}
\alpha\\
\beta
\end{pmatrix}.
\end{equation}
The upper equation is automatic, while the lower one gives \(-\beta=\beta\), hence \(\beta=0\). Normalization requires \(|\alpha|=1\); choosing the irrelevant overall phase so that \(\alpha=1\),
\begin{equation}
|+z\rangle=
\begin{pmatrix}
1\\
0
\end{pmatrix}.
\end{equation}

For eigenvalue \(-1\),
\begin{equation}
\begin{pmatrix}
\alpha\\
-\beta
\end{pmatrix}
=
\begin{pmatrix}
-\alpha\\
-\beta
\end{pmatrix}.
\end{equation}
Now the lower equation is automatic and the upper one gives \(\alpha=-\alpha\), so \(\alpha=0\). Choosing \(\beta=1\) gives
\begin{equation}
|-z\rangle=
\begin{pmatrix}
0\\
1
\end{pmatrix}.
\end{equation}

Now change the experiment and measure \(S_x\). Since
\begin{equation}
\sigma_x
\begin{pmatrix}
\alpha\\
\beta
\end{pmatrix}
=
\begin{pmatrix}
0\cdot\alpha+\beta\\
\alpha+0\cdot\beta
\end{pmatrix}
=
\begin{pmatrix}
\beta\\
\alpha
\end{pmatrix},
\end{equation}
the \(+1\) equation becomes
\begin{equation}
\begin{pmatrix}
\beta\\
\alpha
\end{pmatrix}
=
\begin{pmatrix}
\alpha\\
\beta
\end{pmatrix}.
\end{equation}
The two component equations, \(\beta=\alpha\) and \(\alpha=\beta\), are the same condition. Normalization gives
\begin{equation}
2|\alpha|^2=1,
\end{equation}
so a convenient choice is
\begin{equation}
|+x\rangle
=
\frac{1}{\sqrt2}
\begin{pmatrix}
1\\
1
\end{pmatrix}.
\end{equation}

\begin{figure}[t]
\centering
\includegraphics[width=0.74\textwidth]{lecture_08_figure_03.png}
\caption{For eigenvalue \(+1\), \(\sigma_x\) exchanges the two components, requiring \(\alpha=\beta\).}
\lectureframe{00:39:42}
\end{figure}

\begin{classroomqa}
\lecturetimestamp{00:40:13}
\audiencequestion{If the factor of one-half were included, would the eigenvalue be one-half?}
\lecturerresponse{Yes. The Pauli matrix \(\sigma_x\) has eigenvalues \(\pm1\), while \(S_x=\sigma_x/2\) has eigenvalues \(\pm\tfrac12\).}
\end{classroomqa}

For eigenvalue \(-1\),
\begin{equation}
\begin{pmatrix}
\beta\\
\alpha
\end{pmatrix}
=
-
\begin{pmatrix}
\alpha\\
\beta
\end{pmatrix},
\end{equation}
so
\begin{equation}
\beta=-\alpha,
\qquad
\alpha=-\beta.
\end{equation}
These are again the same equation. After normalization,
\begin{equation}
|-x\rangle
=
\frac{1}{\sqrt2}
\begin{pmatrix}
1\\
-1
\end{pmatrix}.
\end{equation}
Thus an equal superposition of \(z\)-up and \(z\)-down is polarized along the positive \(x\)-axis, while the same superposition with a relative minus sign is polarized along the negative \(x\)-axis.

For the \(y\)-direction,
\begin{equation}
\sigma_y
\begin{pmatrix}
\alpha\\
\beta
\end{pmatrix}
=
\begin{pmatrix}
-i\beta\\
i\alpha
\end{pmatrix}.
\end{equation}
The two normalized eigenvectors are
\begin{equation}
|+y\rangle
=
\frac{1}{\sqrt2}
\begin{pmatrix}
1\\
i
\end{pmatrix},
\qquad
|-y\rangle
=
\frac{1}{\sqrt2}
\begin{pmatrix}
1\\
-i
\end{pmatrix},
\end{equation}
up to an overall phase. Both states give probability \(1/2\) for \(z\)-up and \(1/2\) for \(z\)-down. The relative phase \(i\), which is invisible in those two probabilities, distinguishes \(y\)-polarization from \(x\)-polarization. Linear combinations also describe spin polarized along arbitrary directions; that calculation is deferred.

\begin{classroomqa}
\lecturetimestamp{00:44:22}
\audiencequestion{Has the sign of the \(i\) component been assigned to the wrong \(\sigma_y\) eigenvalue?}
\lecturerresponse{The sign may indeed have been assigned to the wrong eigenvalue. The assignment is uncertain here; it is not being fixed from memory.}
\end{classroomqa}

The direct check is
\begin{align}
\sigma_y\frac{1}{\sqrt2}
\begin{pmatrix}
1\\
i
\end{pmatrix}
&=
\frac{1}{\sqrt2}
\begin{pmatrix}
-i\,i\\
i
\end{pmatrix}
=
\frac{1}{\sqrt2}
\begin{pmatrix}
1\\
i
\end{pmatrix},\\
\sigma_y\frac{1}{\sqrt2}
\begin{pmatrix}
1\\
-i
\end{pmatrix}
&=
\frac{1}{\sqrt2}
\begin{pmatrix}
(-i)(-i)\\
i
\end{pmatrix}
=
-\frac{1}{\sqrt2}
\begin{pmatrix}
1\\
-i
\end{pmatrix}.
\end{align}
Thus \(+i\) belongs to eigenvalue \(+1\), and \(-i\) belongs to eigenvalue \(-1\).\footnote{The spoken assignment remains unsettled after several reversals; the signs shown here are fixed by the displayed matrix multiplication.} Every spin-one-half state, including every state polarized along an arbitrary direction, is a linear combination of \(|+z\rangle\) and \(|-z\rangle\).

\section{Spin One and Position-Dependent Spinors}

For spin \(1\), the possible \(S_z\) values are \(+1,0,-1\). The state space is therefore three-dimensional. Spin \(1\) is integer spin and corresponds to bosonic statistics. We seek three \(3\times3\) matrices satisfying the same angular-momentum commutators.

A representation symmetric among the three coordinate directions is given below.\footnote{The sign of \(S_y\) follows the matrix visible near \(00{:}48{:}50\) and is also fixed by \([S_x,S_y]=iS_z\). Reversing that sign would reverse the required commutator.}
\begin{equation}
S_x
=
i
\begin{pmatrix}
0&0&0\\
0&0&1\\
0&-1&0
\end{pmatrix},
\qquad
S_y
=
i
\begin{pmatrix}
0&0&1\\
0&0&0\\
-1&0&0
\end{pmatrix},
\qquad
S_z
=
i
\begin{pmatrix}
0&1&0\\
-1&0&0\\
0&0&0
\end{pmatrix}.
\end{equation}
Each matrix is \(i\) times a real antisymmetric matrix. Each contains one \(1\), one \(-1\), and a vanishing row and column associated with its own coordinate: the first for \(S_x\), the second for \(S_y\), and the third for \(S_z\). This gives a useful mnemonic. In this representation \(S_z\) is not diagonal.

Let
\begin{equation}
v=
\begin{pmatrix}
\alpha\\
\beta\\
\gamma
\end{pmatrix}.
\end{equation}
Then
\begin{equation}
S_zv
=
\begin{pmatrix}
i\beta\\
-i\alpha\\
0
\end{pmatrix}.
\end{equation}
For eigenvalue \(+1\),
\begin{equation}
i\beta=\alpha,
\qquad
-i\alpha=\beta,
\qquad
0=\gamma.
\end{equation}
Thus \(\beta=-i\alpha\) and \(\gamma=0\). Normalization gives \(2|\alpha|^2=1\), so
\begin{equation}
|m=+1\rangle
=
\frac{1}{\sqrt2}
\begin{pmatrix}
1\\
-i\\
0
\end{pmatrix}.
\end{equation}
For eigenvalue \(-1\),
\begin{equation}
i\beta=-\alpha,
\qquad
-i\alpha=-\beta,
\qquad
0=-\gamma,
\end{equation}
hence \(\beta=i\alpha\), \(\gamma=0\), and
\begin{equation}
|m=-1\rangle
=
\frac{1}{\sqrt2}
\begin{pmatrix}
1\\
i\\
0
\end{pmatrix}.
\end{equation}
Finally, the zero-eigenvalue equation gives \(\alpha=\beta=0\), leaving \(\gamma\) arbitrary. Its normalized solution is
\begin{equation}
|m=0\rangle
=
\begin{pmatrix}
0\\
0\\
1
\end{pmatrix}.
\end{equation}
There are no further eigenvalues: these three eigenvectors exhaust the three-dimensional space. The factors \(1/\sqrt2\) are not cosmetic. The same coefficients later enter quantitative experimental predictions when the algebra is applied to isospin.

\begin{classroomqa}
\lecturetimestamp{00:54:10}
\audiencequestion{How should the complex entries of these spin-one eigenvectors be read?}
\lecturerresponse{For \(m=+1\) the components are \(1/\sqrt2,-i/\sqrt2,0\); for \(m=-1\) they are \(1/\sqrt2,i/\sqrt2,0\); and for \(m=0\) they are \(0,0,1\).}
\end{classroomqa}

Other triplets of matrices satisfy the same algebra and are physically equivalent. The symmetric representation above is useful because the three directions are treated alike, and it will be used again.

\subsection{Restoring position}

Now restore the particle's motion. Spin commutes with position and with momentum, so either may be specified together with a spin component. Position and momentum do not commute, and distinct spin components do not commute with one another.

A moving spin-one-half particle is described by a position-dependent spinor,
\begin{equation}
\psi(x)=
\begin{pmatrix}
\psi_\uparrow(x)\\
\psi_\downarrow(x)
\end{pmatrix}.
\end{equation}
The quantities
\begin{equation}
\rho_{\uparrow_z}(x)=|\psi_\uparrow(x)|^2,
\qquad
\rho_{\downarrow_z}(x)=|\psi_\downarrow(x)|^2
\end{equation}
are joint position-and-spin probability densities. Their sum,
\begin{equation}
\rho(x)=\psi^\dagger(x)\psi(x)
=|\psi_\uparrow(x)|^2+|\psi_\downarrow(x)|^2,
\end{equation}
is the position density. If the particle is known to have been detected at \(x\), then the conditional probabilities are
\begin{equation}
P(\uparrow_z\mid x)
=
\frac{|\psi_\uparrow(x)|^2}{\psi^\dagger(x)\psi(x)},
\qquad
P(\downarrow_z\mid x)
=
\frac{|\psi_\downarrow(x)|^2}{\psi^\dagger(x)\psi(x)},
\end{equation}
whenever \(\rho(x)\neq0\). Thus an electron is represented by a two-component spatial wavefunction, with the local spin amplitudes allowed to vary from place to place.

For a moving spin-one particle, write
\begin{equation}
\phi(x)=
\begin{pmatrix}
\phi_1(x)\\
\phi_2(x)\\
\phi_3(x)
\end{pmatrix}.
\end{equation}
At every point in space, the three components form a vector in the same spin space used above. Spin-one particles have integer spin and are bosons.

\begin{classroomqa}
\lecturetimestamp{01:00:10}
\audiencequestion{Is \(\phi_1(x)\) itself the probability for the spin-one outcome \(m=+1\)?}
\lecturerresponse{No. It is an amplitude in the chosen three-component basis. To obtain a definite spin outcome, form the bra of the corresponding normalized eigenvector, project the field onto it, and take the modulus squared.}
\end{classroomqa}

\begin{figure}[t]
\centering
\includegraphics[width=0.82\textwidth]{lecture_08_figure_04.png}
\caption{Normalized \(\sigma_y\) eigenvectors alongside the three spin-one-half matrices.}
\lectureframe{01:00:12}
\end{figure}

For example,
\begin{equation}
|m=+1\rangle
=
\frac{1}{\sqrt2}
\begin{pmatrix}
1\\
-i\\
0
\end{pmatrix}
\quad\Longrightarrow\quad
\langle m=+1|
=
\frac{1}{\sqrt2}
\begin{pmatrix}
1&i&0
\end{pmatrix}.
\end{equation}
The projection amplitude is therefore
\begin{align}
\langle m=+1|\phi(x)\rangle
&=
\frac{1}{\sqrt2}
\begin{pmatrix}
1&i&0
\end{pmatrix}
\begin{pmatrix}
\phi_1(x)\\
\phi_2(x)\\
\phi_3(x)
\end{pmatrix}\\
&=
\frac{\phi_1(x)+i\phi_2(x)}{\sqrt2}.
\end{align}
Applying the same rule to all three eigenvectors gives the spin-resolved densities.\footnote{These formulas complete the projection algebra after the explicit correction that \(\phi_1\) and \(\phi_2\) need not be real. Treating them as real would incorrectly erase the interference that distinguishes the \(m=+1\) and \(m=-1\) densities.}
\begin{align}
\rho_{+1}(x)
&=
\frac12\left|\phi_1(x)+i\phi_2(x)\right|^2,\\
\rho_{-1}(x)
&=
\frac12\left|\phi_1(x)-i\phi_2(x)\right|^2,\\
\rho_0(x)
&=
|\phi_3(x)|^2.
\end{align}
They satisfy
\begin{equation}
\rho_{+1}(x)+\rho_{-1}(x)+\rho_0(x)
=
|\phi_1(x)|^2+|\phi_2(x)|^2+|\phi_3(x)|^2
=
\phi^\dagger(x)\phi(x).
\end{equation}
If the particle is conditioned to have been detected at \(x\), then
\begin{equation}
P(m\mid x)
=
\frac{\rho_m(x)}{\phi^\dagger(x)\phi(x)},
\end{equation}
provided \(\phi^\dagger(x)\phi(x)\neq0\). The general measurement rule is always the same: form the bra of the desired normalized eigenvector, project the actual state onto it, and take the modulus squared. The resulting density retains its position dependence.

Higher-spin matrices are more complicated, but the principle does not change.

\begin{classroomqa}
\lecturetimestamp{01:07:31}
\audiencequestion{How large is the state space for spin three-halves?}
\lecturerresponse{It has four states, with \(m=\tfrac32,\tfrac12,-\tfrac12,-\tfrac32\). Spin two has five states.}
\end{classroomqa}

\begin{classroomqa}
\lecturetimestamp{01:07:55}
\audiencequestion{Why has the discussion involved \(x,y,z\) but no time direction? Is this treatment nonrelativistic?}
\lecturerresponse{So far it is nonrelativistic. Spin, like mass, is most simply characterized in the particle's rest frame; its description in another inertial frame is obtained by the appropriate Lorentz transformation. That relativistic step is deferred.}
\end{classroomqa}

The Dirac equation supplies the relativistic setting and explains why its four-component wavefunction contains a spin doublet.

\section{Dirac Matrices and the Four-Component Spinor}

The four-component Dirac wavefunction satisfies
\begin{equation}
i\frac{\partial\Psi}{\partial t}
=
-i\alpha_i\frac{\partial\Psi}{\partial x_i}
+m\beta\Psi,
\qquad i=1,2,3,
\end{equation}
where a repeated spatial index is summed. Thus the derivative term means
\begin{equation}
\alpha_1\frac{\partial\Psi}{\partial x}
+\alpha_2\frac{\partial\Psi}{\partial y}
+\alpha_3\frac{\partial\Psi}{\partial z}.
\end{equation}
For a plane wave this becomes the algebraic equation
\begin{equation}
\omega\Psi=(\alpha_i k_i+m\beta)\Psi.
\end{equation}

\begin{classroomqa}
\lecturetimestamp{01:11:32}
\audiencequestion{Do the two occurrences of \(m\) denote the same thing?}
\lecturerresponse{No. One was being used as a spatial index and the other is the mass of the particle. Rename the spatial index \(i\); when \(i\) occurs twice it is summed from \(1\) to \(3\).}
\end{classroomqa}

Squaring the frequency operator must reproduce the relativistic relation
\begin{equation}
\omega^2=\mathbf{k}^2+m^2.
\end{equation}
The square terms therefore require, for each \(i\),
\begin{equation}
\alpha_i^2=I_4,
\qquad
\beta^2=I_4.
\end{equation}

\begin{classroomqa}
\lecturetimestamp{01:13:01}
\audiencequestion{Is \(\beta\) different for each of the four components of \(\Psi\)?}
\lecturerresponse{No. There is one matrix \(\beta\), acting on the whole four-component spinor, and it satisfies \(\beta^2=I_4\).}
\end{classroomqa}

The mixed terms must vanish. For different spatial directions, and also between every \(\alpha_i\) and \(\beta\), this gives
\begin{equation}
\alpha_i\alpha_j+\alpha_j\alpha_i=0
\quad (i\neq j),
\qquad
\alpha_i\beta+\beta\alpha_i=0.
\end{equation}
Equivalently,
\begin{equation}
\{\alpha_i,\alpha_j\}=2\delta_{ij}I_4,
\qquad
\{\alpha_i,\beta\}=0,
\qquad
\{\beta,\beta\}=2I_4.
\end{equation}
Each matrix anticommutes with every different matrix, while its anticommutator with itself is twice the identity.

Without \(\beta\), the three Pauli matrices would already provide three \(2\times2\) matrices that square to the identity and anticommute pairwise. The mass term demands a fourth such matrix. There is no fourth \(2\times2\) matrix with the required properties, and \(3\times3\) matrices do not solve the algebra either. The smallest representation is therefore \(4\times4\).

A \(4\times4\) matrix may be divided into four \(2\times2\) blocks. In one convenient representation,
\begin{equation}
\alpha_i=
\begin{pmatrix}
\sigma_i&0\\
0&-\sigma_i
\end{pmatrix},
\qquad i=1,2,3.
\end{equation}

\begin{figure}[t]
\centering
\includegraphics[width=0.68\textwidth]{lecture_08_figure_05.png}
\caption{The Dirac matrix \(\alpha_i\) written in \(2\times2\) blocks, with \(\sigma_i\) and \(-\sigma_i\) on the diagonal.}
\lectureframe{01:17:23}
\end{figure}

\begin{classroomqa}
\lecturetimestamp{01:17:41}
\audiencequestion{Were the \(\sigma_i\) blocks arranged the other way in the earlier representation?}
\lecturerresponse{Yes. They were placed off the diagonal in another representation. The two arrangements are equivalent and are related by a similarity transformation.}
\end{classroomqa}

In the present representation the remaining matrix is
\begin{equation}
\beta=
\begin{pmatrix}
0&I_2\\
I_2&0
\end{pmatrix},
\qquad
I_2=
\begin{pmatrix}
1&0\\
0&1
\end{pmatrix}.
\end{equation}
Block multiplication reduces checks of the \(4\times4\) algebra to ordinary multiplication of \(2\times2\) matrices.

\begin{classroomqa}
\lecturetimestamp{01:19:17}
\audiencequestion{Did the earlier representation use a different sign or a diagonal form for \(\beta\)?}
\lecturerresponse{It used a different but physically equivalent basis. A basis change mixes the four components of \(\Psi\) and changes the appearance of the matrices without changing the physics.}
\end{classroomqa}

The four entries of \(\Psi\) are not the four coordinates of spacetime. They are coordinates in the smallest vector space on which this anticommutation algebra can act. There are two distinct twofold structures: the division into \(2\times2\) blocks and the two-component Pauli structure within each block. Anticipating the interpretation established below, the four possibilities may be organized as
\begin{equation}
\begin{array}{c|cc}
& S_z=+\tfrac12 & S_z=-\tfrac12\\ \hline
\omega>0 & (\omega>0,+\tfrac12) & (\omega>0,-\tfrac12)\\
\omega<0 & (\omega<0,+\tfrac12) & (\omega<0,-\tfrac12)
\end{array}.
\end{equation}
The lower row consists of negative-frequency electron states. It is not a row of positrons; a positron will be interpreted as a hole in the filled set of those states.

\section{Rest-Frame Energy Doubling and the Positron}

The twofold energy structure is easiest to expose at zero momentum. When \(\mathbf{k}=0\), the plane-wave state is spatially constant and the spatial-derivative term disappears. The rest-frame equation\footnote{The frame at \(01{:}23{:}38\) retains a transient extra factor of \(i\) on the mass term; the spoken correction at \(01{:}24{:}05\) removes it.} is
\begin{equation}
i\frac{\partial\Psi}{\partial t}=m\beta\Psi.
\end{equation}

\begin{figure}[t]
\centering
\includegraphics[width=0.76\textwidth]{lecture_08_figure_06.png}
\caption{At zero momentum the spatial derivatives vanish, leaving the corrected rest-frame equation \(i\partial_t\Psi=m\beta\Psi\).}
\lectureframe{01:23:38}
\end{figure}

Divide the four-component column into an upper and a lower two-component object:
\begin{equation}
\Psi=
\begin{pmatrix}
\psi_+\\
\psi_-
\end{pmatrix}.
\end{equation}

\begin{classroomqa}
\lecturetimestamp{01:24:47}
\audiencequestion{Do the labels \(\psi_+\) and \(\psi_-\) already mean positive and negative energy?}
\lecturerresponse{No. At this stage they are only names for the upper and lower two-component blocks. Their energy combinations have not yet been found.}
\end{classroomqa}

\begin{classroomqa}
\lecturetimestamp{01:25:10}
\audiencequestion{What does the divided four-component column represent?}
\lecturerresponse{It is the same four-component Dirac spinor, organized to match the block form of the matrices. Each half is a two-component object, and two such objects make the original four-component column.}
\end{classroomqa}

The matrix \(\beta\) swaps the two blocks:
\begin{equation}
\beta
\begin{pmatrix}
\psi_+\\
\psi_-
\end{pmatrix}
=
\begin{pmatrix}
\psi_-\\
\psi_+
\end{pmatrix}.
\end{equation}
The rest-frame equation therefore becomes
\begin{equation}
i\frac{\partial\psi_+}{\partial t}=m\psi_-,
\qquad
i\frac{\partial\psi_-}{\partial t}=m\psi_+.
\end{equation}

There are two natural suggestions. Taking another time derivative would decouple the system, but adding and subtracting the two equations diagonalizes them immediately. Indeed,
\begin{align}
i\frac{\partial}{\partial t}(\psi_++\psi_-)
&=m(\psi_++\psi_-),\\
i\frac{\partial}{\partial t}(\psi_+-\psi_-)
&=-m(\psi_+-\psi_-).
\end{align}
Defining
\begin{equation}
\chi_+=\psi_++\psi_-,
\qquad
\chi_-=\psi_+-\psi_-,
\end{equation}
gives two independent equations,
\begin{equation}
i\partial_t\chi_+=m\chi_+,
\qquad
i\partial_t\chi_-=-m\chi_-.
\end{equation}
For definite-frequency dependence \(e^{-i\omega t}\), the sum has
\begin{equation}
\omega=+m,
\end{equation}
while the difference has
\begin{equation}
\omega=-m.
\end{equation}
This is exactly the positive- and negative-energy doubling found in the earlier one-dimensional example: the relativistic equation again has positive-energy and negative-energy electron solutions. Thus the block doubling separates into positive- and negative-energy branches. Pure rest-frame states may be written
\begin{equation}
\Psi_{\omega=+m}=
\begin{pmatrix}
\xi\\
\xi
\end{pmatrix},
\qquad
\Psi_{\omega=-m}=
\begin{pmatrix}
\xi\\
-\xi
\end{pmatrix},
\end{equation}
where \(\xi\) is still an arbitrary two-component object.

The sum and difference have disposed of one twofold structure, but the two components of \(\xi\) remain. The Pauli matrices act on this residual pair. This is the pair associated with spin, subject to the rotation and conservation-law check discussed below.

\begin{classroomqa}
\lecturetimestamp{01:32:23}
\audiencequestion{Can spin be treated as the outer twofold structure and energy sign as the inner one?}
\lecturerresponse{Yes. The two factor orderings are equivalent. One may interchange which index is written first; the change is only a reshuffling of the four components.}
\end{classroomqa}

The negative-frequency branch leads to the Dirac-sea interpretation. A negative frequency is assigned negative electron energy. Filling every negative-energy electron state lowers the energy as far as the exclusion principle permits, since only one electron can occupy each complete fermionic state. The completely filled negative-energy sea is then called the vacuum.

Promote one electron from this filled sea into a positive-energy state. The result contains a positive-energy electron and an unoccupied negative-energy state. Removing a negative energy raises the energy, so the hole has positive excitation energy; because an electron is missing, it also carries the opposite electric charge. The hole is a positron.

\begin{classroomqa}
\lecturetimestamp{01:32:57}
\audiencequestion{Can the positron interpretation be explained once more?}
\lecturerresponse{The sum of the two blocks has positive frequency and the difference has negative frequency. The negative-frequency solutions are negative-energy electron states, and the vacuum is obtained by filling all of them. Promoting one filled electron to positive energy leaves both a positive-energy electron and a hole. That positive-energy, oppositely charged hole is the positron.}
\end{classroomqa}

The algebra has therefore produced two consequences at once: positive- and negative-energy branches, which become the particle--antiparticle structure, and a residual two-component structure associated with spin one-half. The striking point is that both emerged from the smallest matrices capable of satisfying the relativistic anticommutation relations. Spin had only recently appeared as an empirical property when this mathematical construction produced particles with spin together with an antiparticle branch; in that historical setting the result was dramatic.

\begin{classroomqa}
\lecturetimestamp{01:35:14}
\audiencequestion{Are the negative-energy electrons themselves positrons?}
\lecturerresponse{No. The negative-energy electron states are filled in the vacuum. A positron is the absence of one electron from that filled set: a hole, not a negative-energy electron.}
\end{classroomqa}

A hole may participate in transitions. A positron has positive energy because it is the absence of a negative-energy electron. If a positive-energy electron fills the hole, the electron and positron annihilate, but their energy must be carried away, ordinarily by photons. In sea language, an electron has filled an unoccupied negative-energy state. The same release of energy can be pictured as an excited electron in an atom dropping to a lower level: the downward transition emits photons.

\begin{classroomqa}
\lecturetimestamp{01:36:39}
\audiencequestion{Why do higher negative-energy electrons not fall into the hole?}
\lecturerresponse{That question mixes two cases. A positive-energy electron filling the hole is electron--positron annihilation, with photons carrying away the energy. A transition by a negative-energy sea electron is a different process, taken up below.}
\end{classroomqa}

The zero-momentum calculation by itself has established a twofold degeneracy after energy sign is separated. A degeneracy of two states is not, by itself, a proof that those states carry angular momentum.

\begin{classroomqa}
\lecturetimestamp{01:37:44}
\audiencequestion{How does the zero-momentum calculation prove that the remaining pair of states is spin?}
\lecturerresponse{The pair may be called up and down, and the Pauli matrices appear in the full equation, but that alone is not yet the proof. One must show either how the pair transforms under spatial rotations or that assigning it spin makes total angular momentum conserved in interactions and collisions. That demonstration is deferred.}
\end{classroomqa}

For a positive-energy rest state the upper and lower blocks are equal,
\begin{equation}
\Psi_{\omega=+m}=
\begin{pmatrix}
\xi\\
\xi
\end{pmatrix}.
\end{equation}
Establishing that \(\xi\) is an angular-momentum doublet requires the rotation or conservation argument just described; it does not follow merely from counting two components.

\begin{classroomqa}
\lecturetimestamp{01:40:46}
\audiencequestion{Can a different electron annihilate with the positron even if its energy is different?}
\lecturerresponse{Yes. Electrons retain no individual identity. Either available electron may fill the hole, and differences in the initial energy and momentum are reflected in the number and energies of the emitted photons.}
\end{classroomqa}

The hole picture also explains a transition wholly within the negative-energy sea. Suppose an occupied negative-energy state moves into the hole. The original hole is filled, but the state from which that electron came is now empty. The hole has moved. Equivalently, the positron has changed its energy and has emitted radiation.

\begin{classroomqa}
\lecturetimestamp{01:42:12}
\audiencequestion{Why cannot neighboring negative-energy levels successively fall into the hole?}
\lecturerresponse{They can move the hole from one level to another. An electron dropping into the old hole leaves a new hole at its former energy, so the initial positron becomes a positron of different energy. The energy difference is carried by photons.}
\end{classroomqa}

\begin{classroomqa}
\lecturetimestamp{01:44:39}
\audiencequestion{Why does this not produce an endless downward radiative cascade?}
\lecturerresponse{Because energy and momentum must be conserved simultaneously. Those two conservation laws do not permit the unrestricted process. The same issue arises for an ordinary free electron: it cannot simply keep lowering its energy by emitting photons in empty space.}
\end{classroomqa}
